\section{学习方法和题解复盘的深层要求}

刷题教材不能只给最终代码。读者真正需要的是从题目约束到算法选择的推理链：输入规模允许什么复杂度，数据是否有序，值域能否压缩，状态是否会重复，局部选择能否证明，容器操作成本是否符合题目热路径。每道题都要先写暴力解，因为暴力解暴露了最朴素的状态空间；然后再指出哪个查询、哪个枚举、哪个重复子问题最贵。优化不是把模板贴上去，而是把最贵的那一步替换为更合适的数据结构或更强的不变量。

高质量复盘至少包含三层。第一层是题面复述，用自己的话说明输入、输出和约束，尤其要把隐藏条件翻译成算法语言，例如正数数组带来窗口单调性，排序数组带来二分和双指针，树节点唯一性决定能否用值建索引。第二层是复杂度预算，把 n、m、值域、字符串长度、图边数分开，不要只写一个 n。第三层是正确性说明，用不变量、交换论证、归纳或最短路层次证明解释为什么不会漏答案。

题解里的代码要服务于解释，而不是炫技。变量名应该表达状态含义，例如 left 表示窗口左边界，farthest 表示当前可达最远位置，indegree 表示尚未满足的依赖数量。复杂代码要先拆出检查函数、松弛函数或合并函数。这样做不是为了形式，而是让读者能在面试中逐句讲清楚代码对应的数学状态。代码写完以后，至少用一个最小用例、一个边界用例、一个反例用例和一个规模用例手算。

三个月计划只能作为节奏，不应该成为内容边界。真正的学习路径是循环式的：第一轮建立题型地图，第二轮补齐证明和边界，第三轮把相似题横向比较，第四轮把代码改成工程上也能接受的版本。每一轮都要记录失败原因。失败原因不是笼统的不会，而是要拆成读错题、复杂度预算错误、状态定义不完整、边界初始化错误、容器 API 误用、证明缺失或代码实现失误。

面试表达要像讲工程方案。先说暴力方案，说明它为什么慢；再说关键观察，说明它让哪一步变便宜；然后给出数据结构职责，说明保存什么、查询什么、删除什么；最后给出不变量和复杂度。这个顺序能防止答案变成模板堆砌。面试官追问时，优先回到约束和不变量，而不是机械背另一段代码。

实验是把知识落地的关键。数组题打印指针区间，窗口题打印计数表和 formed，二分题打印 left、mid、right 和谓词结果，图题打印队列层次，DP 题打印小表格，堆题打印堆顶和过期状态。只要一次实验能解释一个常见错误，它就不是额外负担，而是把题解变成长期能力的工具。

\begin{principlebox}{专题复盘要求}
读完本节后，不要只记结论。请至少回到 LeetCode 53 最大子数组和、LeetCode 560 和为 K 的子数组、LeetCode 875 爱吃香蕉、LeetCode 312 戳气球这类代表题，把每个结论还原成一个可验证的问题：它解决了哪一步重复工作，依赖哪个题目条件，使用哪个容器或状态，边界条件如何破坏错误写法。
\end{principlebox}
